% The robotic sixth finger on the affected wrist; redrawn from paper/figures/tikz with the
% report's type sizes.
\begin{tikzpicture}[scale=0.62, line width=0.8pt]
  \draw[fill=black!5] (-1.6,-4) -- (-1.4,-1) -- (1.6,-1) -- (1.8,-4);
  \draw[fill=black!5, rounded corners=8pt] (-2,-1.2) rectangle (2,2.2);
  \foreach \x/\h in {-1.5/2.6, -0.5/3.1, 0.5/2.9, 1.5/2.3} {
    \draw[fill=black!5, rounded corners=5pt] (\x-0.38,2.0) rectangle (\x+0.38,2.0+\h);
    \draw[black!30, line width=0.4pt] (\x-0.38,2.0+\h*0.45) -- (\x+0.38,2.0+\h*0.45);
    \draw[black!30, line width=0.4pt] (\x-0.38,2.0+\h*0.75) -- (\x+0.38,2.0+\h*0.75);
  }
  \draw[fill=black!5, rounded corners=5pt, rotate around={40:(-2,0.2)}] (-2.4,0.2) rectangle (-1.6,2.6);
  \draw[fill=black!30] (-1.6,-1.9) rectangle (1.8,-1.4);
  \draw[draw=rred, fill=rred!15, rounded corners=3pt] (2.0,-1.8) rectangle (2.7,0.6);
  \draw[draw=rred, fill=rred!15, rounded corners=3pt] (2.05,0.8) rectangle (2.65,2.4);
  \draw[draw=rred, fill=rred!15, rounded corners=3pt] (2.1,2.6) rectangle (2.6,3.7);
  \foreach \y in {0.7,2.5} \fill[rred] (2.35,\y) circle (0.12);
  \node[font=\small, text=rred, anchor=west, align=left] at (3.1,1.3) {robotic sixth finger,\\on the affected wrist};
  \node[font=\small, anchor=west] at (3.1,-1.65) {wrist cuff};
\end{tikzpicture}
